\section{Conclusion and Limitations}
\label{sec:conclusion}

Low-bit VLA deployment is a policy-configuration problem: the mask changes the policy, and the policy changes the observations it meets. We formulated this as byte-constrained policy configuration under a mask-induced rollout distribution and instantiated a proposal--admission separation in \method: RIPA proposes precision from representation integrity, FCP projects the layout to an exact byte budget and admits complete candidates only on full-policy temporal evidence, and \dyrange resolves the remaining closed-loop activation-range shift with a frozen mask. On the shared RoboCasa365 evaluation, the method attains the highest observed success among the evaluated PTQ methods on both backbones: 54.0\% at 2.22$\times$ audited Linear-component compression on GR00T and 27.7\% at 2.70$\times$ on $\pi_{0.5}$. With the mask and quantized weights held fixed, input-conditioned A8 improves success by 10.6 points over a pooled static-A8 calibration.

\paragraph{Limitations.}
Fusion and projection-rule superiority remain statistically unresolved. Controlled probes verify the temporal sensitivities that D-PAC targets, while the matched-candidate replay selects the same anchor under all three criteria. The calibration inputs and frozen-bank allocation are reproduced; the historical candidate-ranking stage has not been regenerated, and complete offline configuration costs are unavailable. Evaluation is simulation-only, some baselines require non-equivalent adapters or different measurement scopes, and the generic backend trades memory for latency and energy. The demonstrated benefit is component-storage compression at competitive success; the whole-checkpoint ratio is smaller (Appendix~\ref{sec:storage_scope}).
\label{sec:main-paper-end}

% Self-review: method-led, sources retained, component claims bounded by evidence.
